\chapter{The Single Electron at the Pairing Door}

The recoil route closed \(\alpha\) by letting an atom answer with motion. The next use of the apparatus asks a
different kind of question. Bring the named electron to a superconducting boundary. On one side the ledger reads
a single charged excitation: odd parity, one electron's worth of charge, a tanged disturbance against the even
background. On the other side the ledger reads a condensate of Cooper pairs: charge moving in units of \(2e\),
phase coherent, even parity, a many-body state whose individual electrons are no longer separately addressed.

The arrow
\[
  e^- \longrightarrow e^- \longrightarrow \hbox{Cooper pair}
\]
must be read carefully. It does not say that one electron literally turns into two electrons by itself. It says
that the manual now names the second electron at the junction before it reads the superconducting ledger. An
outer electron can meet a junction electron; that junction electron can then meet the pair boundary. The two
electron positions both read as electron in the finite gauge, but they are not the same stage of the run. The
apparatus must record which stage it is reading and what residual the boundary paid to make the reading
consistent.

\section{What Has to Be Added}

Three new slots are needed. The first is parity. A Cooper-pair condensate is even: charges enter its coherent
ledger two electron charges at a time. A lone electron is odd. The difference is not cosmetic. It is the bit that
quasiparticle poisoning, parity switching, and pair-breaking experiments all watch. The apparatus should not
fold this bit into a smooth amplitude. It should tange it.

The second slot is phase. A superconducting pair state carries a phase, and a pair transfer changes the phase
account by one turn in the simplified ledger. The phase is not the same thing as the count of pairs. They are
conjugate faces: count asks how many pair transfers have been recorded; phase asks where the coherent condensate
stands. The manual model keeps only a small integer phase-turn counter, enough to show the bookkeeping seam
without pretending to derive BCS theory.

The third slot is cost. The electron-pair boundary is cyclic. If the apparatus is allowed to ask forever, it can
move forever around the same four-stage trace: outer electron, junction electron, Cooper pair, returning
electron. That is not a bug in the model. It is the point. A superconducting contact is not a one-shot lookup
table; it is an interface that can be driven repeatedly. The only honest finite computation is therefore a
budgeted one: decide how much cost to spend, run that many folds, and report where the interface stopped.

\section{The Lean Face}

The companion file is \(\texttt{device/Measurement/CooperManual.lean}\). It begins with three faces:

\begin{quote}\small
\begin{verbatim}
inductive Face where
  | singleElectron
  | junctionElectron
  | cooperPair
deriving Repr, DecidableEq
\end{verbatim}
\end{quote}

The faces are then read through the existing Agent/Positron gauge:

\begin{quote}\small
\begin{verbatim}
def Face.reading : Face -> Reading
  | .singleElectron => electron
  | .junctionElectron => electron
  | .cooperPair => superconductingPair
\end{verbatim}
\end{quote}

This is a deliberate reuse of the apparatus already present in the Lean files. The electron is the all-funge
corner of the gauge. The superconducting pair is the reading one tange off that corner, at the magnitude gate.
The interface manual now does invent a second named electron position: the junction electron. It is not a new
particle species in the gauge. It is the second electron slot in the interaction trace. Both electron faces read
as \(\texttt{electron}\), while the \(\texttt{Stage}\) record remembers whether the run is at the outer electron
or at the junction.

The two counters are:

\begin{quote}\small
\begin{verbatim}
def Face.positronCount (face : Face) : Nat :=
  Measurement.Agent.positronCount electron face.reading

def Face.matterCount (face : Face) : Nat :=
  Measurement.Agent.matterReading electron face.reading
\end{verbatim}
\end{quote}

The name \(\texttt{positronCount}\) is inherited from the existing gauge vocabulary. In this interface it should
be read more narrowly as the tange count: the number of gates that no longer funge with the electron reading.
Both electron faces have no tange in the current apparatus. The Cooper-pair face has exactly one. That one tange
is the admitted oddness of the pair boundary, the place where the electron-reading side and the even condensate
fail to be the same reading.

\section{What Lean Does Not Supply}

Lean does not supply the superconducting gap \(\Delta\), the Josephson energy \(E_J\), the charging energy
\(E_C\), or the gate offset \(n_g\). It does not decide whether the physical device is a Josephson junction, a
Cooper-pair box, a transmon, a nanowire, or an Andreev contact. Those are laboratory choices. The Lean file only
gives the interface a shape into which those choices can be placed.

The shape is useful precisely because it is small. If the laboratory wants pair tunneling, the pair face records
the transfer. If the laboratory wants parity poisoning, an electron face records the odd event. If the laboratory
wants an electron-electron contact before the Cooper-pair read, the junction electron carries that residual. If
the laboratory wants to compare repeated switching, the run is budgeted by cost. In all cases the first
discipline is the same: say which stage is active, say what count moved, say what phase turn was paid, and never
hide the fact that the interface can keep going.
